\section{Bemessung nach Stommel et al.}
\label{sec:stommel_ergebnis}

\subsection{Methode wie veröffentlicht}
\label{sec:stommel_veroeffentlicht}

Der Nachweis nach \secref{sec:auswertung_beanspruchung} wurde auf alle 99 Stellungen je Serie des Taschenkörpers aus \secref{sec:waelzweg} angewendet. Die Lastgeschichte jedes Knotens besteht aus dem unbelasteten Zustand und den vollständigen Stellungen eines Lastschritts. Eingestellt waren der Überlastungsfall F2, der zyklische Weg über die Hauptspannungen und die Abminderungsklasse teilkristallin verstärkt für das Verfahren wie veröffentlicht. Berichtet werden die beiden Fußrundungen des mittleren Zahns T3 mit 2562 Knoten. Die Randzähne tragen den Sektorrand. \tabref{tab:stommel_veroeffentlicht} nennt je Lastschritt den Knoten der größten zyklischen Auslastung $a_{\mathrm{B,V}}$ mit seiner statischen Auslastung $a_{\mathrm{stat}}$ und der vorhergesagten Lastwechselzahl auf der synthetischen Wöhlerlinie. W2 ist die Hauptlinie, W1 die lineare Referenz. Daneben steht die gemessene Lastwechselzahl der Bruchlinie. Sie ist keine Anrisslinie, der Vergleich ist deshalb einseitig und wird in \secref{sec:genauigkeit} in drei Stufen geführt.

\begin{table}[htbp]
    \caption{Nachweis nach Stommel et al. wie veröffentlicht an den Fußrundungen des Zahns T3, quasistatische Karte bei $80$\,\textdegree C, Überlastungsfall F2: kritischer Knoten mit Lage, zyklische Auslastung $a_{\mathrm{B,V}}$, statische Auslastung $a_{\mathrm{stat}}$ dort, Zahl der Fußknoten über $1$ und vorhergesagte Lastwechselzahl $N_{\mathrm{pred}}$. Daneben die gemessene Lastwechselzahl $N_{\mathrm{L},50\%}$ der Bruchlinie aus \secref{sec:woehlerlinien}.}
    \label{tab:stommel_veroeffentlicht}
    \centering
    \footnotesize
    \setlength{\tabcolsep}{4pt}
    \begin{tabular}{L{0.8cm}L{1.0cm}L{2.9cm}L{1.1cm}L{1.1cm}L{1.5cm}L{1.5cm}L{1.5cm}L{1.8cm}}
        Stufe & $T_{2}$ in Nm & kritischer Knoten        & $a_{\mathrm{B,V}}$ & $a_{\mathrm{stat}}$ & Fußknoten $a_{\mathrm{B,V}} > 1$ & Fußknoten $a_{\mathrm{stat}} > 1$ & $N_{\mathrm{pred}}$  & $N_{\mathrm{L},50\%}$ gemessen \\
        \hline
        \arrayrulecolor{black!30}
        W1    & $8$           & Fuß T2/T3, $z = 0{,}00$  & $1{,}72$           & $0{,}84$            & $1536$                           & $0$                               & $1{,}6 \cdot 10^{3}$ & $2{,}67 \cdot 10^{7}$          \\
        \hline
        W1    & $12$          & Fuß T3/T4, $z = 0{,}00$  & $2{,}69$           & $1{,}06$            & $2312$                           & $176$                             & sofort               & $5{,}96 \cdot 10^{6}$          \\
        \hline
        W1    & $16$          & Fuß T3/T4, $z = 0{,}00$  & $10{,}40$          & $1{,}64$            & $2486$                           & $944$                             & sofort               & $1{,}36 \cdot 10^{6}$          \\
        \hline
        W2    & $8$           & Fuß T3/T4, $z = -5{,}93$ & $1{,}30$           & $0{,}42$            & $905$                            & $0$                               & $1{,}7 \cdot 10^{5}$ & $2{,}67 \cdot 10^{7}$          \\
        \hline
        W2    & $12$          & Fuß T3/T4, $z = -5{,}93$ & $1{,}64$           & $0{,}52$            & $1943$                           & $0$                               & $4{,}3 \cdot 10^{3}$ & $5{,}96 \cdot 10^{6}$          \\
        \hline
        W2    & $16$          & Fuß T3/T4, $z = -5{,}93$ & $1{,}96$           & $0{,}66$            & $2228$                           & $0$                               & sofort               & $1{,}36 \cdot 10^{6}$          \\
        \hline
        \arrayrulecolor{black}
    \end{tabular}
\end{table}

% TK erledigt (nein, Bruchlinie; Vergleich einseitig in drei Stufen, Satz vor der Tabelle und in 7.1): Bei der Bruchlinie in Abschnitt 6.3 ist da auch tatsächlich mit den ersten Anrissen gerechnet? Sonst kann man schlecht gegen NL 50 % gemessen vergelichen, wenn das Buch und wir in der Erklärung sagen, dass das Verfahren auf ersten Anriss und nicht Bruchmechanik ausgelegt ist und jetzt aber verglichen wird mit der Buchlinie? 

Mit der plastischen isotropen Stufe W2 besteht der statische Nachweis an allen Fußknoten. Die zyklische Auslastung liegt bei $1{,}30$, $1{,}64$ und $1{,}96$. Die vorhergesagte Lastwechselzahl fällt mit dem Moment von $1{,}7 \cdot 10^{5}$ über $4{,}3 \cdot 10^{3}$ auf sofortiges Versagen bei $16$\,Nm, wo sie die Grenze von $10^{4}$ Lastspielen des Verfahrens unterschreitet. Die lineare Stufe W1 verletzt bei $12$ und $16$\,Nm an 176 und 944 Fußknoten den statischen Nachweis und sagt schon bei $8$\,Nm nur $1{,}6 \cdot 10^{3}$ Lastwechsel vorher. Der kritische Knoten von W2 liegt in der Fußrundung zwischen T3 und T4 nahe der Stirnseite bei $z = -5{,}93$\,mm von $\pm 7{,}5$\,mm. Er sitzt am Fuß der belasteten Flanke des Zahns T4 mit $h = 0{,}49$. Unter dem eigenen Zahn steht er im Zug, unter dem Nachbarzahn im Druck. Dass der Knoten zu T4 gehört, folgt aus seinem Winkel von $94{,}1$° um die Radachse, $2{,}9$° von der Zahnmitte T4 und $4{,}1$° von der Zahnmitte T3 entfernt. Dazu kommt die Kontaktspur des Modells, die auf der Flanke jedes Zahns zur Seite des kleineren Winkels liegt. Bei W1 liegt er in der Zahnmitte. Auf dem Ringkörper der Voruntersuchungen lag der kritische Knoten von W2 auf der Druckseite derselben Fußrundung mit $R \approx -0{,}4$ und einer Vorhersage von $3{,}5 \cdot 10^{5}$ bis $1{,}7 \cdot 10^{3}$ Lastwechseln. Der Steg des Taschenkörpers verschiebt die Lastgeschichte der Fußknoten so, dass die Zugseite maßgebend wird.

\begin{figure}[htbp]
    \centering
    \includegraphics[width=\textwidth]{kap5/stommel_karte_8nm_w2}
    \caption{Zyklische Auslastung $a_{\mathrm{B,V}}$ über der Oberfläche zwischen den Mitten der Zähne T3 und T4, also über der Fußrundung, den beiden zugewandten Flanken und den Kopfflächen, Winkel um die Radachse über der Zahnbreite, W2 elastisch-plastisch isotrop, quasistatisch, $8$\,Nm. Der leere Streifen zwischen Kopffläche und Flanke bei $90{,}7$ bis $90{,}9$° und $96{,}0$ bis $96{,}2$° ist die Kopffase der STplus-Kontur, die das Netz mit einer Elementreihe ohne Zwischenknoten überspannt.}
    \label{fig:stommel_karte_8nm}
\end{figure}

% TK erledigt (Karte neu: Grid hinten, Zahnmitten T3/T4 als Linien, Grenzlinie a_B,V = 1 mit Legende, Ausschnitt auf die Zahnmitten; Schwarzweiß bleibt eine Farbkarte, die Grenzlinie trägt die Aussage): Schlecht, wenn schwarz weiß, hast dud a eine gute Idee? Wenn nicht, dann passt das schön, sieht gut aus. Aber Grid bitte auch weider in den Hintergrund und kannst du noch senksrechte Linien einfügen mit den Bezeichnungen T3 und T4, also direkt bei der mittleren Position des Zahnkopfes, also Zahnmitte? Dann kann man sich noch besser hineinversetzen, da der Leser erstmal an sich nicht weiß, wo jetzt der 0 Grad Zahn sein soll und den Winkel um die Radachse so nicht nullen kann im Kopf. Verstehst du?

\begin{figure}[htbp]
    \centering
    \includegraphics[width=\textwidth]{kap5/stommel_verlauf_8nm_w2}
    \caption{Vergleichsspannung $\sigma_{\mathrm{V}}$ und vorzeichenbehaftete von-Mises-Spannung der beiden kritischen Knoten der größten zyklischen und der größten statischen Auslastung über dem Wälzweg, W2, quasistatisch, $8$\,Nm. Senkrechte Linien an den Normpunkten A bis E.}
    \label{fig:stommel_verlauf_8nm}
\end{figure}

\figref{fig:stommel_karte_8nm} trägt die zyklische Auslastung von W2 bei $8$\,Nm über der Fußrundung zwischen T3 und T4 und über den angrenzenden Flanken bis zu den Zahnmitten auf. Die Auslastung über eins liegt in zwei Bändern entlang der Fußrundung, an 955 der 83\,650 Knoten. Das Band bei $94{,}1$° gehört zum Fuß der belasteten Flanke des Zahns T4. Es trägt den Größtwert $1{,}30$ in zwei stirnseitennahen Knoten bei $z = \pm 5{,}93$\,mm auf dem Radius $24{,}78$\,mm. Das Band bei $92{,}8$° gehört zum Fuß der unbelasteten Flanke des Zahns T3, der im Druck steht, solange T3 trägt, und in den Zug kommt, sobald T4 trägt. Im Band bei $94{,}1$° fällt die Auslastung zur Zahnmitte hin auf $1{,}29$, und die Linie $a_{\mathrm{B,V}} = 1$ schließt das Band gegen die Druckseite und gegen die Flanke ab. \figref{fig:stommel_verlauf_8nm} zeigt die Lastgeschichte der beiden kritischen Knoten über dem Wälzweg. Der Knoten der größten zyklischen Auslastung liegt in dieser Fußrundung. Er steht im Zug, solange der Zahn T4 trägt, mit $33$\,MPa Spitze, und wechselt in den Druck, sobald die Last auf den Nachbarzahn T3 übergeht. Daraus folgen die Amplitude $28{,}3$\,MPa und die Mittelspannung $4{,}6$\,MPa. Der Knoten der größten statischen Auslastung ist ein Kopfknoten des Zahns T1 mit $50$\,MPa Vergleichsspannung unter dem Kontakt des Ritzels. Seine Spannung tritt in einer Stellung auf und bleibt mit $0{,}77$ unter der statischen Grenze.

Gegen die Bruchlinie des Laufversuchs liegt die Vorhersage mit W2 bei $8$\,Nm um den Faktor $160$ und bei $12$\,Nm um den Faktor $1400$ zu niedrig, bei $16$\,Nm sagt sie sofortiges Versagen vorher. Das Verhältnis wächst mit dem Moment. Der Nachweis endet beim ersten makroskopischen Anriss, der Laufversuch beim Bruch. Die Anrissfenster aus \secref{sec:woehlerlinien} verkleinern den Abstand bei $8$\,Nm auf den Faktor $20$ bis $46$. Bei $12$ und $16$\,Nm wurden die ersten Anrisse nicht ausdrücklich dokumentiert, so dass dort nur eine obere Schranke bleibt und die Prüfanweisung in \chapref{chap:zusammenfassung} die Dokumentation der Anrisse für die folgenden Läufe verlangt. Die Einordnung dieser Abstände folgt in \secref{sec:genauigkeit}.

% TK erledigt (Satz zur fehlenden Dokumentation bei 12 und 16 Nm mit Verweis auf die Prüfanweisung): Ah, hier hast du es schon erklärt. Vielleicht noch darauf eingehen kurz, dass die Anrisse für 12 und 16 Nm leider nicht ausdrücklich dokumentiert wurden und dies für die folgenden Versuche zwingend dokumentiert werden sollte. sieeh ausblick oder wo du es eh schon dann mit erwähnt hast.

\subsection{Anisotrope Erweiterung}
\label{sec:stommel_anisotrop}

Für die anisotropen Stufen W3 und W4 wurde die Vergleichsspannung durch die Hill-Vergleichs\-spannung im Faserrahmen mit allen sechs Verhältnissen der Karte ersetzt, nach \secref{sec:auswertung_beanspruchung}. Die Orientierung ist in allen Nachweisen dieses Abschnitts die des CONVERSE-Mappings, nicht die der CT-Zonen. Die Abminderungsfaktoren nehmen die Klasse teilkristallin unverstärkt, weil die Karte an Zugversuchen des verstärkten Werkstoffs kalibriert ist und die Klasse der verstärkten Werkstoffe die Querschwäche ein zweites Mal zählen würde. \tabref{tab:stommel_hill} nennt die Ergebnisse an den Fußrundungen des Zahns T3 auf dem Taschenkörper, dazu W1 unter Hill als Methodenvergleich mit denselben Spannungen wie in \tabref{tab:stommel_veroeffentlicht}.

\begin{table}[htbp]
    \caption{Nachweis mit der Hill-Vergleichs\-spannung im Faserrahmen an den Fußrundungen des Zahns T3, Klasse teilkristallin unverstärkt, quasistatisch, Überlastungsfall F2: kritischer Knoten, zyklische und statische Auslastung, Fußknoten über $1$. $\infty$ bedeutet eine Vergleichsspannung jenseits des Endes der Isotherme bei $80$\,\textdegree C.}
    \label{tab:stommel_hill}
    \centering
    \footnotesize
    \setlength{\tabcolsep}{4pt}
    \begin{tabular}{L{1.6cm}L{1.2cm}L{3.2cm}L{1.3cm}L{1.3cm}L{1.8cm}L{1.8cm}L{1.6cm}}
        Stufe   & $T_{2}$ in Nm & kritischer Knoten        & $a_{\mathrm{B,V}}$ & $a_{\mathrm{stat}}$ & Fußknoten $a_{\mathrm{B,V}} > 1$ & Fußknoten $a_{\mathrm{stat}} > 1$ & $N_{\mathrm{pred}}$ \\
        \hline
        \arrayrulecolor{black!30}
        W3 Hill & $8$           & Fuß T3/T4, $z = -0{,}67$ & $3{,}25$           & $\infty$            & $2390$                           & $54$                              & sofort              \\
        \hline
        W3 Hill & $12$          & Fuß T3/T4, $z = 1{,}33$  & $4{,}67$           & $\infty$            & $2505$                           & $614$                             & sofort              \\
        \hline
        W3 Hill & $16$          & Fuß T3/T4, $z = 1{,}33$  & $55{,}1$           & $\infty$            & $2544$                           & $1116$                            & sofort              \\
        \hline
        W4 Hill & $8$           & Fuß T2/T3, $z = 2{,}79$  & $2{,}45$           & $0{,}45$            & $2128$                           & $0$                               & sofort              \\
        \hline
        W4 Hill & $12$          & Fuß T2/T3, $z = 2{,}79$  & $3{,}01$           & $0{,}66$            & $2411$                           & $0$                               & sofort              \\
        \hline
        W4 Hill & $16$          & Fuß T2/T3, $z = 2{,}79$  & $3{,}55$           & $\infty$            & $2498$                           & $9$                               & sofort              \\
        \hline
        W1 Hill & $8$           & Fuß T2/T3, $z = 0{,}67$  & $3{,}29$           & $\infty$            & $2271$                           & $112$                             & sofort              \\
        \hline
        W1 Hill & $12$          & Fuß T3/T4, $z = 2{,}00$  & $5{,}68$           & $\infty$            & $2467$                           & $588$                             & sofort              \\
        \hline
        W1 Hill & $16$          & Fuß T3/T4, $z = -2{,}79$ & $\infty$           & $\infty$            & $2546$                           & $1049$                            & sofort              \\
        \hline
        \arrayrulecolor{black}
    \end{tabular}
\end{table}

Unter Hill wandert der kritische Knoten von der Stirnseite in die Zahnmitte, bei W3 nach $z = -0{,}7$ und $1{,}3$\,mm, bei W4 nach $2{,}8$\,mm und bei W1 nach $0{,}7$ bis $2{,}8$\,mm. Dort liegen nach \secref{sec:faserorientierung} die Fasern der Flankenzonen axial, die Biegespannung in der Stirnebene wirkt quer zur Faser, und die Karte schwächt diese Richtung. An der Stirnseite folgen die Fasern der Kontur und tragen die Biegung in Faserrichtung. Dass der kritische Knoten ohne Hill nahe der Stirnseite liegt, hat seinen Grund im Spannungszustand. An der freien Stirnfläche verschwindet die Spannung normal zur Stirnseite. Der zweiachsige Zustand liefert bei ähnlicher größter Hauptspannung eine höhere von-Mises-Spannung und einen kleineren Mehraxialitätsgrad als der dreiachsige Zustand in der Zahnmitte, in dem die Querdehnung behindert ist. Die Hill-Vergleichs\-spannung liegt an den Fußknoten deutlich über der von-Mises-Spannung. Die Fließspannungsverhältnisse der Karte betragen im Mittel über die Knoten $R_{11} = 0{,}88$, $R_{22} = 0{,}69$, $R_{33} = 0{,}57$, $R_{12} = 0{,}58$, $R_{13} = 0{,}54$ und $R_{23} = 0{,}46$, so dass die Schubspannung im Faserrahmen fast doppelt gewichtet wird. Bei W3 und bei W1 unter Hill überschreitet die Vergleichsspannung an 54 bis 1116 Fußknoten das Ende der Isotherme, so dass dort keine Dehnung abgelesen werden kann und der statische Nachweis entfällt. Die plastische anisotrope Stufe W4 besteht den statischen Nachweis bis $12$\,Nm, bei $16$\,Nm liegen neun Fußknoten jenseits der Isotherme. Zyklisch liegt sie mit $2{,}5$ bis $3{,}5$ um den Faktor $1{,}8$ über W2. \figref{fig:stommel_verlauf_16nm_w4} zeigt den Verlauf der Hill-Vergleichsspannung der beiden über das ganze Rad kritischen Knoten von W4 bei $16$\,Nm, der Knoten der größten zyklischen und der größten statischen Auslastung. Dazu ist die vorzeichenbehaftete von-Mises-Spannung mit dem Vorzeichen der betragsgrößten Hauptspannung aufgetragen. Beide sind Kopfknoten unter dem Ritzelkontakt an den Stirnseiten der Zähne T1 und T5, nicht die Fußknoten der Tabelle, die der Nachweis unter Hill jenseits der Isotherme findet. Der Knoten der größten statischen Auslastung liegt auf der Druckseite, seine Kurve ist deshalb negativ. Die Höhe der Auslastung hängt an der Klasse der Abminderungsfaktoren, die das Buch für einen Werkstoff mit aufgelöster Richtungsabhängigkeit nicht festlegt. Mit der Klasse der verstärkten Werkstoffe läge sie noch höher. Belastbar sind damit zunächst die Lage des kritischen Knotens und der Vergleich der Stufen. Um die Erweiterung auch für die Lebensdauer zu kalibrieren, braucht es die Orte der ersten Anrisse über der Zahnbreite aus den Laufversuchen. Die Prüfanweisung in \chapref{chap:zusammenfassung} verlangt sie.

% TK erledigt (Kurven im Text erklärt: Hill-Vergleichsspannung beider kritischer Knoten plus vorzeichenbehaftete von-Mises-Spannung, Druckknoten negativ; letzter Satz auf die Kalibrierung über die Anrissorte gerichtet): Also bild 5.9 ist für mich noch nicht eindeutig, vorallem die Bezeichnungen Vergelichsspannung und hill. Mises ist doch die vergelichsspannung im normalfall oder nicht? Und warum ist die eine so positiv und die andere so negativ? Und der letzte Satz passt auch nicht in die Richtung, die ich untersuche. Mit dem Satz würde ich so hinnehmen, dass das Verfahren nicht für die Lebensdauer verwendbar ist, aber genau das versuchen wir ja hinzubekommen und jetzt ist ja eh noch mit faktoren gerechnet, die noch nie für sowas angewendet wurden. Aber da bräuchte es mehr Versuchsdaten, um die Anrisse besser zu lokalisieren und auf Faktoren zurückzuschließen oder?


% TK erledigt (Stirnseite: ebener Spannungszustand, höhere von-Mises-Spannung und kleineres h; Mitte unter Hill: Fasern axial, Biegung quer zur Faser): Warum ist ohne Hill überhaupt der höchste Auslastungsgrad eher bei der Stirnseite, das kann ich mir nicht wirklich erklären, sind die Effekte da durch Mehraxialität und anderen Einflüssen auf den dynamischen Auslastungsgrad so prägsam, dass die größeren Spannungen in der Zahnmitte trotzdem nicht so relevant sind? Bitte mal erklären, warum sich das dann unter Hill auch so verschiebt. Achso, bei Hill, da die mittlere Zone die schwächste ist. Ok, aber warum dann bei mises eher Stirnseite, das checke ich noch nicht?

\begin{figure}[htbp]
    \centering
    \includegraphics[width=\textwidth]{kap5/stommel_verlauf_16nm_w4hill}
    \caption{Hill-Vergleichs\-spannung und vorzeichenbehaftete von-Mises-Spannung der beiden über das ganze Rad kritischen Knoten (größte zyklische und größte statische Auslastung, beide Kopfknoten unter dem Ritzelkontakt) über dem Wälzweg, W4 elastisch-plastisch anisotrop, Taschenkörper, quasistatisch, $16$\,Nm. Senkrechte Linien an den Normpunkten A bis E.}
    \label{fig:stommel_verlauf_16nm_w4}
\end{figure}

\subsection{Globalmaxima, Überlastungsfälle und Vergleich mit S-Life Plastics}
\label{sec:stommel_global_slife}

Über alle 83\,650 Knoten des Rades liegt die größte Auslastung bei W2 in der Fußrundung selbst, mit $1{,}96$ bei $16$\,Nm. Die Kopfknoten der Zähne auf dem Radius $25{,}80$\,mm, wo die gerundete Ritzelkopfkante die Radflanke belastet, folgen mit $1{,}93$ am Kopf von T2 dicht dahinter. Bei W1 geht $a_{\mathrm{B,V}}$ an diesen Kopfknoten gegen unendlich. Unter Hill kommt die erste Knotenreihe unter dieser Spur auf $25{,}78$\,mm hinzu, mit Vergleichsspannungen jenseits der Isotherme. Die Kopfknoten sind Kontaktspannungszustände und keine Aussage zur Zahnfußermüdung. Die Tabellen dieses Abschnitts beziehen sich deshalb auf die Fußrundungen des Zahns T3 und nennen das Globalmaximum daneben. Im Versuch ist genau diese Spur der Ort der Einarbeitung der Ritzelkopfkante und bei $16$\,Nm von vier Flankenbrüchen an drei Rädern (\secref{sec:schadensbilder}). Der Überlastungsfall F2 ist nach \secref{sec:auswertung_beanspruchung} der Lastpfad des Rades mit konstantem Spannungsverhältnis $R$ an jedem Knoten. Die plastischen Stufen weichen von der Proportionalität zwischen Moment und Spannung leicht ab.

% TK erledigt (Schadensbilder in 6.2 aufgenommen und hier sowie in 7.2 mit der Kantenspur verknüpft; 5.3 rechnet mit des CONVERSE-Mappings, so im Text benannt): Da wäre jetzt doch interessant, mal die Schadensbilder 70_input\71_Messungen\Mikroskop\Bilder\Schadensbilder_Kunststoffraeder\8Nm\F4092_1_Lw_170824_202610.jpg 70_input\71_Messungen\Mikroskop\Bilder\Schadensbilder_Kunststoffraeder\12Nm\F4092_1_Lw_170817_202610.jpg 70_input\71_Messungen\Mikroskop\Bilder\Schadensbilder_Kunststoffraeder\16Nm\F4092_1_Lw_170820_202610.jpg zu analysieren, da das genau der Bereich sein könnte, in dem sich das Phänomen auch bei den realen Laufrädern ereignet und sich die Ritzelkopfkante in das Kunststoffrad einarbeitet. Was schon top wäre für die Auswertung. Ist 5.3.3 schon mit dem CT Orientierungsprofil oder dem eigenen Spritzgegossenem? Aber das ist natürlich interessant. Bitte da nochmal reinblicken und das dann mit aufnehmen in den späteren Auswertungskapiteln.

Als Gegenprobe wurde die Stellung C mit W1 in S-Life Plastics mit den Einstellungen nach \secref{sec:auswertung_beanspruchung} ausgewertet, rein schwellend mit $R = 0$, die Spannungen ungemittelt als Elementknotenwert eines Elements, die kritische Ebene mit $11{,}25$° Schritt. Verglichen wurden vier Knoten, zwei Zugknoten der Fußrundung in Zahnmitte und nahe der Stirnseite, ein Druckknoten der gegenüberliegenden Fußrundung und ein Kopfknoten. \tabref{tab:slife_vergleich} nennt die Werte bei $16$\,Nm. Die Hauptspannungen stimmen bis auf die Mittelung überein, mit $82{,}84$ gegen $82{,}85$\,MPa am Zugknoten der Zahnmitte und bis $5$\,\% am Kopfknoten. Der statische Nachweis stimmt auf $5$\,\%. Die verbleibende Differenz stammt aus der Dehnung, die S-Life auf seiner angepassten Isotherme liest, $2{,}49$ gegen $2{,}62$\,\% der interpolierten Punkte der Karte. Am Kopfknoten unter Druck mit $h = -1{,}3$ gibt S-Life $1{,}01$ gegen $0{,}42$ des eigenen Nachweises, der die Druckhauptspannungen nach Gleichung (6.17) durch $1{,}3$ teilt und die zulässige Dehnung bei $h = -1{,}3$ liest. S-Life wendet dort ein anderes Kriterium an. Im zyklischen Nachweis sind Amplitude, Wechselfestigkeit und Mittelspannungsfaktor gleich, mit $41{,}41$ gegen $41{,}42$\,MPa, $20{,}12$ gegen $20{,}19$\,MPa und $0{,}82$ gegen $0{,}815$. Der einzige Unterschied ist die statische Kappe der zulässigen Amplitude. Die beiden Lesarten unterscheiden sich in der Größe, von der die Mittelspannung abgezogen wird. S-Life nimmt die statische Zulässigkeit des unbelasteten Extremzustands von $47{,}32$\,MPa, also die zulässige Spannung der Isotherme ohne Rücksicht auf den Spannungszustand des Knotens, und erhält am Zugknoten $5{,}91$\,MPa. Der eigene Nachweis nimmt die statische Bauteilfestigkeit $\sigma_{\mathrm{stat}}(h)$ des maßgebenden Zustands, die das Buch im statischen Nachweis mit dem Mehraxialitätsgrad des Knotens liest, und erhält $13{,}27$\,MPa. Daraus folgen $a_{\mathrm{B,V}}$ von $6{,}53$ gegen $2{,}67$. Auf der Druckseite übersteigt die Mittelspannung von $49{,}6$\,MPa die Kappe von S-Life, so dass dessen Auslastung ins Unendliche läuft, während der eigene Nachweis $10$\,MPa behält. Die Umsetzung ist damit in Spannungen, statischem Nachweis, Amplitude und Festigkeit bestätigt. Der Unterschied in der Kappe ist benannt und reproduzierbar. Das Buch legt die Lesart in den Gleichungen (6.24) bis (6.26) \cite[S.~374]{Sto18} nicht eindeutig fest. Einen Maßstab für die Genauigkeit nennt es an anderer Stelle, und zwar für den Vergleich mit dem Versuch, nicht für den Unterschied zweier Lesarten. Für die Ermüdungsfestigkeit gilt eine Abweichung der berechneten von der gemessenen Lastspielzahl um den Faktor zwei bis drei als gut und das Treffen der Größenordnung als angemessen \cite[S.~375]{Sto18}. An diesem Maßstab gemessen liegt die vorhergesagte Lastwechselzahl der eigenen Lesart nach \tabref{tab:stommel_veroeffentlicht} bei $8$ und $12$\,Nm zwei bis drei Größenordnungen unter der Bruchlinie, also außerhalb des Bandes. Die Lesart von S-Life mit der kleineren Kappe läge noch weiter darunter. Der Unterschied der Lesarten, Faktor $2{,}5$ in der Auslastung und über die Neigung $k$ von rund $15$ mehrere Größenordnungen in der Lastspielzahl, ändert an dieser Einordnung nichts, bleibt aber als Unterschied benannt. Die Arbeit führt deshalb den eigenen Nachweis als wörtliche Lesart des Buches und S-Life als Referenzimplementierung nebeneinander, ohne eine der beiden auszuschließen.
% TK erledigt (Lesart der Kappe im Text erklärt, S-Life gegen eigene Lesart, Einordnung über die Genauigkeitsaussage des Buchs S. 375 wie vom Nutzer zitiert): Meinst du, welches der sigma a,i ... was du für den Nachweis verwendet hast und was die verwenden? Aber in der folgenden Seite steht auch: "Speziell für den Bereich der Ermüdungsfestigkeit ... Faktoren von 2 bis 3 ... als gut zu bezeichnen sind." Also können

\begin{table}[htbp]
    \caption{S-Life Plastics 4.0.0 gegen den eigenen Nachweis an vier Knoten der Stellung C, W1, quasistatische Karte, $16$\,Nm. Je Zelle erster Wert S-Life, zweiter Wert eigener Nachweis: größte Hauptspannung, statische Auslastung, Amplitude, Wechselfestigkeit, angewandte zulässige Amplitude und zyklische Auslastung. S-Life mit ungemittelten Elementwerten, eigener Nachweis mit gemittelten Elementknotenwerten.}
    \label{tab:slife_vergleich}
    \centering
    \footnotesize
    \setlength{\tabcolsep}{4pt}
    \begin{tabular}{L{2.6cm}L{1.9cm}L{1.6cm}L{1.9cm}L{1.8cm}L{2.0cm}L{1.6cm}}
        Knoten                      & $\sigma_{\mathrm{1}}$ bzw. $\sigma_{\mathrm{3}}$ in MPa & $a_{\mathrm{stat}}$ & $\sigma_{\mathrm{a}}$ in MPa & $\sigma_{\mathrm{W}}$ in MPa & $\sigma_{\mathrm{B}}$ angewandt in MPa & $a_{\mathrm{B,V}}$  \\
        (S-Life / eigener Nachweis) &                                                         &                     &                              &                              &                                        &                     \\
        \hline
        \arrayrulecolor{black!30}
        Zug, Zahnmitte              & $82{,}84$ / $82{,}85$                                   & $1{,}39$ / $1{,}45$ & $41{,}41$ / $41{,}42$        & $20{,}12$ / $20{,}19$        & $5{,}91$ / $13{,}27$                   & $6{,}53$ / $2{,}67$ \\
        \hline
        Zug, nahe Stirnseite        & $84{,}19$ / $81{,}67$                                   & $1{,}42$ / $1{,}49$ & $42{,}06$ / $40{,}84$        & $20{,}12$ / $20{,}19$        & $5{,}26$ / $14{,}80$                   & $7{,}52$ / $2{,}42$ \\
        \hline
        Druck, Fußrundung           & $-99{,}12$ / $-99{,}17$                                 & $1{,}08$ / $1{,}11$ & -- / $3{,}62$                & $20{,}12$ / $20{,}19$        & $0$ / $26{,}12$                        & $\infty$ / $4{,}60$ \\
        \hline
        Kopf                        & $-79{,}88$ / $-76{,}29$                                 & $1{,}01$ / $0{,}42$ & -- / $14{,}47$               & $20{,}12$ / $20{,}19$        & -- / $26{,}12$                         & $4{,}23$ / $1{,}07$ \\
        \hline
        \arrayrulecolor{black}
    \end{tabular}
\end{table}

% TK erledigt (Bildunterschrift und Kopfzeile nennen die Reihenfolge S-Life / eigener Nachweis): Also bei der Tabelle 5.5 sollte schon erkennbar sein auf den ersten Blick, was eigenes verfahern und S-Life ist. Man kann es sich mit dem Text deuten, aber die Tabelle sollete auch für sich alleine sprechen. Noch in der Tabellen beschriftung aufnehmen, erster wert, zweiter wert oder innerhalb der Tabelle anders kennzeichen.s

% Herkunft:
% 5.3.1: analyse.md Abschnitt 7 (Tabelle je Serie: W1 8/12/16 Nm 1,63/2,03/4,47, a_stat 0,81/0,94/1,31, Knoten > 1
%   1438/2295/2479 und 0/256/920, N 3,91e3/0/0, kritische Knoten 13999 (T2_T3, z −5,25), 5817 (T3_T4, z −6,75);
%   W2 1,24/1,47/1,74, 0,48/0,58/0,68, 660/1489/2097, 0/0/0, 3,54e5/2,24e4/1,65e3, Knoten 10852/10853 (T3_T4,
%   z −6,00)); ADR-118 Amendment 2026-10-12 Punkt 2 (Einstellungen: F2, Hauptspannungsroute, Klasse verstärkt),
%   Punkt 4 (Druckseite, R < 0, z −6,0 bis −6,75), Amendment 2026-10-13 (R ≈ −0,4: +32/−14 MPa bei 8 Nm);
%   N_L,50 % aus 80_output/messungen/woehler/woehler.md (2,671e7, 5,964e6, 1,357e6); Faktoren 2,671e7/3,54e5 =
%   75, 5,964e6/2,24e4 = 266, 1,357e6/1,65e3 = 822; Anrissfenster 8 Nm: 3,22e6/3,54e5 = 9,1 bis 7,66e6/3,54e5
%   = 21,6 (kst-E_69), 3,65e6/3,54e5 = 10,3 bis 6,61e6/3,54e5 = 18,7 (kst-E_75).
% 5.3.2: analyse.md Abschnitt 7 (W3 Hill 3,12/3,85/10,38 mit ∞, Knoten 37549 (T2_T3, z −0,75), 21518 (T3_T4, z
%   −3,75), > 1: 2326/2518/2560 und 74/530/1138; W4 Hill 2,16/2,58/3,03, 0,44/0,53/0,73, 54026 (z 2,25), 26552
%   (z −3,00), 1954/2358/2495 und 0/0/4, N 1,22e3/0/0; W1 Hill 3,07/3,96/16,64 mit ∞, 2173/2469/2543 und
%   79/570/1054); ADR-118 Amendment E7 (σ_Hill = 1,74 σ_Mises im Median an C, σ_13 ≈ 37 MPa mit R_13 0,53; Klasse
%   teilkristallin unverstärkt, kein doppeltes Zählen, Buch S. 346); W4 gegen W2: 2,16/1,24 = 1,74, 2,58/1,47 =
%   1,76, 3,03/1,74 = 1,74.
% 5.3.3: ADR-118 Amendment 2026-10-12 Punkt 3 (Kopfknoten r 25,80, a_B,V → ∞ für W1, W2 16 Nm 1,84 T2 head;
%   Hill: Knotenreihe r 25,78, σ_V 93 bis 105 MPa; FE-22), Amendment 2026-10-13 (Überlastfall als Lastpfad, F2,
%   E2 offen); S-Life: 80_output/fe/kst_e/frozen_mesh_20layers_bore33_gerechnet/stommel/vergleich_slife.md 16 Nm
%   (41472: 82,84/82,85, a_SK 1,388/1,452, σ_a 41,41/41,42, σ_WK 20,12/20,19, Kappe 5,91/13,27, a_BK 6,53/2,671;
%   6147: 84,19/81,67, 1,420/1,489, 42,06/40,84, 5,26/14,80, 7,52/2,416; 41143: −99,12/−99,17, 1,079/1,114, σ_a
%   eigen 3,62, Kappe 0 → a_BK 49 641 (als ∞ geschrieben), eigen 26,12/4,602; 5741: −79,88/−76,29, 1,010/0,425,
%   14,47, 26,12, 4,23/1,066); ADR-118 Amendment 2026-10-14 Findings 1 bis 4 (Dehnung 2,49 gegen 2,62 %, K_AK 0,82
%   gegen K_M 0,815, Kappe σ_S,allow,min 47,32 MPa, Gl. (6.17) Faktor 1,3, Lesart Gl. (6.24) bis (6.26) S. 365 bis
%   367 offen, Option cap = min_state).
% Bildwunsch: kap5/stommel_karte_8nm_w2 (stommel_W2_QS/karte_LOAD_8NM_T3_T4.svg), stommel_verlauf_8nm_w2
%   (stommel_W2_QS/verlauf_LOAD_8NM.svg), stommel_verlauf_16nm_w4hill (stommel_hill_W4_QS/verlauf_LOAD_16NM.svg).
% Bibwunsch: Sto18 vorhanden (Gl. (6.17), (6.24) bis (6.26)).
% Nomenklatur: a_B_V, a_stat, N_pred (prüfen: N_pred als vorhergesagte Lastwechselzahl, neu), sigma_a, sigma_W,
%   sigma_B, h, R, R_13 (Hill-Verhältnis, prüfen), sigma_stat, sigma_V vorhanden (3.7).
% Offen: Klassenwahl unter Hill; Überlastungsfälle F1 bis F4 (E2); Lesart der Kappe; die 12-Nm-Knotentabelle des
%   eigenen Nachweises fehlt in der S-Life-Datei (nur Spannungen).
